
In a secret sharing scheme \cite{Shamir79}, a dealer divides a secret into shares such that the secret remains hidden from anyone who holds an ``unauthorized'' set of shares, but can be recovered from any ``authorized'' set of shares.
These properties are referred to as privacy and correctness, respectively.

In our second result, described in \cref{chap:mdss}, we consider a setting where {\em multiple} dealers may simultaneously distribute shares of their respective secrets.
We put forth a notion of {\em multi-dealer secret sharing} (MDSS) that achieves {\em stronger} security and correctness properties in this setting.

We name our security and correctness properties \emph{unlinkability} and \emph{multi-dealer} (MD) \emph{correctness} respectively, and give intuitive descriptions of them here.

\subsection{Unlinkability}
We propose a new security property for secret sharing that is stronger than the standard notion of privacy.
Intuitively, this property requires that given a set of shares from multiple dealers, an adversary cannot associate (i.e., link) any unauthorized subset of shares to a common dealer.

We present two definitions: unlinkability and one-more unlinkability. The first intuitively matches the  guarantee we would like to achieve, while the second is simpler and easier to use.
We will show that the two definitions are equivalent, and therefore that one can use the second without loss of generality.


\subsection{MD-Correctness}
Our correctness definition differs from standard secret sharing in that we must reconstruct secrets given a set of shares from \emph{multiple} dealers. This set may include shares from both sufficient dealers (i.e. those that output enough shares that their secret should be recovered) and \emph{insufficient} dealers (i.e. those that do not).
We require that the reconstruction algorithm outputs exactly the set of secrets shared by the sufficient dealers, except with some failure probability $\mdssCVar$.
Similar to standard secret sharing, we only consider correctness for honestly generated (i.e. via the sharing algorithm) sets of shares.

\subsection{Results}

For our first main result we show that unlinkability is strictly stronger than the standard secret-sharing notion of privacy.

\begin{theorem*}[\cref{thm:mdss-ul-vs-priv} and \cref{thm:mdss-ul-vs-priv-2}, informal]
	Unlinkability is strictly stronger than privacy
\end{theorem*}

Next, we give a construction of MDSS based on the Guruswami-Sudan list-decoder for Reed-Solomon codes \cite{FOCS:GurSud98}.
% Under this construction, unlinkability is broken once a dealer exceeds $\mdssPrivThr$ shares, while correctness is guaranteed at $\mdssRecThr$ shares, where $\mdssRecThr > \mdssPrivThr$.

\begin{theorem*}[\cref{thm:gs-mdss-correctness} and \cref{thm:gs-mdss-ul}, informal]
	There exists an MDSS scheme that recovers all secrets with at least $\mdssRecThr$ shares from a pool of size $\mdssMaxShares$, while keeping all secrets with under $\mdssPrivThr$ shares unlinkable, for all parameters satisfying $\mdssRecThr > \sqrt{\mdssMaxShares \cdot \mdssPrivThr}$ and $\mdssPrivThr + 1 \leq (1 - \epsilon)\mdssRecThr$ for some constant $0 < \epsilon  \leq 1$.
\end{theorem*}

We additionally give a small-fields optimized construction in \cref{sec:mdss-sf} for settings where keeping the size of the shares small is valuable.

\subsection{Related Work}

Both \cite{guillermo2003providing} and \cite{paskin2020cryptographic} study secret sharing schemes that are \emph{anonymous}, i.e.\ those for which shares do not leak the identities of their holders.
However, they do not address the \emph{multi-dealer} setting considered here, where recovery must be performed over a pool of shares that may be of different secrets.

In subsequent work, \cite{C:BGIJL25} generalizes the notion of MDSS to \emph{fully anonymous secret sharing}, and gives constructions for general access structures (i.e.\ those beyond just thresholding).
Additionally, \cite{con2025anonymous} studies the applications of insertion/deletion-resistant Reed-Solomon codes to anonymous secret sharing schemes.
